% =====================================================================
% §1 INTRODUÇÃO
% Arco: contexto/relevância -> escassez/impacto -> problema real ->
%       objetivos + virada mono->bi-objetivo (pivô vs. qualificação)
% =====================================================================
\section{Introdução}

%%% Frame 1.1 — Hook: o setor em números
\begin{frame}{O setor de diagnóstico por imagem em números}

    \begin{block}{Relevância clínica e econômica}
        \begin{itemize}
            \item 2,4 bilhões de exames diagnósticos realizados anualmente \cite{MedicinaSA2024}.
            \item Receita bruta de R\$ 24 bilhões movimentada em 2022 \cite{Abramed_2023}.
            \item Diagnóstico precoce $\rightarrow$ melhor acompanhamento, menor custo relativo \cite{OBRIEN_2007,Yasui_Fujiwara_Okitsu_Yonemitsu_Ito_Yokosuka_2012}.
        \end{itemize}
    \end{block}

    % TODO: falar 30s sobre a escala do setor como gancho.
    \vspace{2.5cm}
\end{frame}

%%% Frame 1.2 — Escassez + impacto operacional
\begin{frame}{Escassez de profissionais e impacto operacional}

    \begingroup
    \setbeamercolor{block title}{fg=white, bg=red!40}
    \setbeamercolor{block body}{bg=red!10}
    \begin{block}{Gargalo de mão de obra especializada \cite{Liebel_Dias_Schneider_SaScielo2021,RSNA_2024}}
        \begin{itemize}
            \item Crescimento estimado de 6\% para 18\% em vagas não preenchidas.
            \item Brasil: $\approx$ 1,51 mil médicos para cada 1,99 milhão de habitantes (2017).
            \item Equipamentos ociosos, filas e pressão sobre margens operacionais \cite{ANS_2024}.
        \end{itemize}
    \end{block}
    \endgroup

    % TODO: conectar escassez -> necessidade de escalas otimizadas.
    \vspace{2cm}
\end{frame}

%%% Frame 1.3 — O problema abordado
\begin{frame}{O problema abordado}

    \begingroup
    \setbeamercolor{block title}{fg=white, bg=orange!40}
    \setbeamercolor{block body}{bg=orange!10}
    \begin{block}{Escalonamento médico em rede de clínicas}
        \begin{itemize}
            \item Múltiplas unidades \textbf{autônomas} que compartilham profissionais.
            \item Escalas feitas \textbf{manualmente} pelo setor de operações: lento, sem garantia de otimalidade.
            \item Relação \textit{vários-para-vários} médico--unidade $\rightarrow$ explosão combinatória.
        \end{itemize}
    \end{block}
    \endgroup

    \begingroup
    \setbeamercolor{block title}{fg=white, bg=cyan!40}
    \setbeamercolor{block body}{bg=cyan!10}
    \begin{block}{Proposta}
        Gerar escalas otimizadas: maximizar atendimento com a equipe existente \emph{e} minimizar remanejamentos entre unidades \cite{Rousseau_Pesant_Gendreau_2002,Robinson_Chen_2003}.
    \end{block}
    \endgroup
\end{frame}

%%% Frame 1.4 — Objetivos + virada mono->bi-objetivo
\begin{frame}{Objetivos da pesquisa}

    \begin{block}{}
        \begin{itemize}
            \item Modelagem matemática do escalonamento médico em múltiplas unidades.
            \item Investigação \textbf{bi-objetivo}: exames não atendidos \emph{vs.} trocas de clínica.
            \item Avaliar três famílias de solução: exata, MIP-heurísticas e meta-heurística.
            \item Contribuir para a literatura de escalonamento médico multiunidade.
        \end{itemize}
    \end{block}

    \begingroup
    \setbeamercolor{block title}{fg=white, bg=UFTblue}
    \setbeamercolor{block body}{bg=UFTblue!10}
    \begin{block}{Evolução desde a qualificação}
        De objetivos tratados isoladamente para uma formulação \textbf{bi-objetiva} resolvida por método $\varepsilon$-restrito.
    \end{block}
    \endgroup
    % TODO: enfatizar que esta é a mudança central do trabalho.
\end{frame}
